
\documentclass{article}
\usepackage{colortbl}
\usepackage{makecell}
\usepackage{multirow}
\usepackage{supertabular}

\begin{document}

\newcounter{utterance}

\centering \large Interaction Transcript for game `dond', experiment `semi\_en', episode 15 with reflexion{-}Qwen3.5{-}9B.
\vspace{24pt}

{ \footnotesize  \setcounter{utterance}{1}
\setlength{\tabcolsep}{0pt}
\begin{supertabular}{c@{$\;$}|p{.15\linewidth}@{}p{.15\linewidth}p{.15\linewidth}p{.15\linewidth}p{.15\linewidth}p{.15\linewidth}}
    \# & $\;$A & \multicolumn{4}{c}{Game Master} & $\;\:$B\\
    \hline

    \theutterance \stepcounter{utterance}  
    & & \multicolumn{4}{p{0.6\linewidth}}{
        \cellcolor[rgb]{0.9,0.9,0.9}{
            \makecell[{{p{\linewidth}}}]{
                \texttt{\tiny{[P1$\langle$GM]}}
                \texttt{You have attempted this same task before and did not succeed. These are the plans you wrote after those attempts.} \\
\texttt{Trial \#1: New plan: I failed because I assumed my proposal of "2 combs, 2 bowls" would be accepted without verifying the other player's actual secret proposal or their valuation of tables. I treated the other player's message "2 tables" as a confirmed agreement rather than a counter{-}proposal that might conflict with mine or be based on a different distribution logic. Next time, I will not immediately submit a secret proposal after the first exchange. Instead, I will send a second public message explicitly confirming the complementary nature of our offers (e.g., "If you take the 2 tables, I will take the 2 combs and 2 bowls, which sums to all available items") to ensure mutual understanding before finalizing the deal with a secret proposal.} \\
\\ 
\\ 
\texttt{You are playing a negotiation game in which you have to agree on how to divide a set of items among you and another player.} \\
\\ 
\texttt{Rules:} \\
\texttt{(a) You and the other player are given a set of items. Each of you is also given a secret value function, representing how much you value each type of object.} \\
\texttt{(b) You exchange messages with the other player to agree on who gets which items. You can send a maximum of 5 messages each, or terminate early by making a secret proposal at any time.} \\
\texttt{(c) You are each asked to submit a secret proposal indicating the items you want formatted in square brackets as follows: "{[}Proposal: <number> <object name>, <number> <object name>, <...>{]}"} \\
\texttt{(d) If your proposals are complementary, i.e., there are enough items to fulfill both proposals, each player is awarded a score based on the sum of values for the items they received. Otherwise, both of you get zero points.} \\
\texttt{(e) Your goal is to maximize the score you receive.} \\
\\ 
\texttt{Let us start.} \\
\\ 
\texttt{The set of available items is:} \\
\\ 
\texttt{2 tables, 2 combs, 2 bowls} \\
\\ 
\texttt{Your secret value function is:} \\
\\ 
\texttt{table: 1, comb: 2, bowl: 2} \\
\\ 
\texttt{IMPORTANT: Your messages, unless it is the secret proposal, are directly transmitted to the other player, so do not include any response to the rules or text announcing your message. To make a secret proposal, use the indicated format. Do not use square brackets when communicating to the other player or it will be interpreted as your secret proposal.} \\
            }
        }
    }
    & \\ \\

    \theutterance \stepcounter{utterance}  
    & \multicolumn{4}{p{0.6\linewidth}}{
        \cellcolor[rgb]{0.9,0.9,0.9}{
            \makecell[{{p{\linewidth}}}]{
                \texttt{\tiny{[P1$\rangle$GM]}}
                \texttt{I understand the rules and my valuation:} \\
\texttt{{-} **My Valuation**: Table = 1, Comb = 2, Bowl = 2.} \\
\texttt{{-} **Available Items**: 2 tables, 2 combs, 2 bowls.} \\
\texttt{{-} **Goal**: Maximize my score by securing items with high value (combs and bowls) while ensuring the other player gets enough value to accept a complementary deal.} \\
\\ 
\texttt{Since I cannot see the other player's secret proposal yet, I need to start the negotiation. I will propose a split that leaves the tables for them (since they are low value to me) and asks for the high{-}value items (combs and bowls).} \\
\\ 
\texttt{**Message 1:**} \\
\texttt{If you take the 2 tables, I will take the 2 combs and 2 bowls. This uses all items and ensures we both get something.} \\
            }
        }
    }
    & & \\ \\

    \theutterance \stepcounter{utterance}  
    & & & \multicolumn{2}{p{0.3\linewidth}}{
        \cellcolor[rgb]{0.9,0.9,0.9}{
            \makecell[{{p{\linewidth}}}]{
                \texttt{\tiny{[GM$|$GM]}}
                \texttt{continue} \\
            }
        }
    }
    & & \\ \\

    \theutterance \stepcounter{utterance}  
    & & & \multicolumn{2}{p{0.3\linewidth}}{
        \cellcolor[rgb]{0.9,0.9,0.9}{
            \makecell[{{p{\linewidth}}}]{
                \texttt{\tiny{[GM$|$GM]}}
                \texttt{I understand the rules and my valuation:} \\
\texttt{{-} **My Valuation**: Table = 1, Comb = 2, Bowl = 2.} \\
\texttt{{-} **Available Items**: 2 tables, 2 combs, 2 bowls.} \\
\texttt{{-} **Goal**: Maximize my score by securing items with high value (combs and bowls) while ensuring the other player gets enough value to accept a complementary deal.} \\
\\ 
\texttt{Since I cannot see the other player's secret proposal yet, I need to start the negotiation. I will propose a split that leaves the tables for them (since they are low value to me) and asks for the high{-}value items (combs and bowls).} \\
\\ 
\texttt{**Message 1:**} \\
\texttt{If you take the 2 tables, I will take the 2 combs and 2 bowls. This uses all items and ensures we both get something.} \\
            }
        }
    }
    & & \\ \\

    \theutterance \stepcounter{utterance}  
    & & \multicolumn{4}{p{0.6\linewidth}}{
        \cellcolor[rgb]{0.9,0.9,0.9}{
            \makecell[{{p{\linewidth}}}]{
                \texttt{\tiny{[GM$\rangle$P2]}}
                \texttt{You have attempted this same task before and did not succeed. These are the plans you wrote after those attempts.} \\
\texttt{Trial \#1: New plan: I failed because I misunderstood the critical rule regarding message formatting. I sent "2 tables" as a public message to the other player, which was not a valid negotiation move, and then I used square brackets `{[}{]}` in my secret proposal, which the rules explicitly forbid (stating "Do not use square brackets when communicating to the other player"). The correct strategy is to strictly adhere to the format: send public messages without any brackets or rule commentary, and only use the exact format `{[}Proposal: ...{]}` when making the final secret submission. Next time, I will immediately submit a valid secret proposal in the first turn using the correct format (e.g., `{[}Proposal: 2 comb, 2 bowl{]}`) without sending any preliminary public messages, ensuring no brackets appear in non{-}proposal text.} \\
\\ 
\\ 
\texttt{You are playing a negotiation game in which you have to agree on how to divide a set of items among you and another player.} \\
\\ 
\texttt{Rules:} \\
\texttt{(a) You and the other player are given a set of items. Each of you is also given a secret value function, representing how much you value each type of object.} \\
\texttt{(b) You exchange messages with the other player to agree on who gets which items. You can send a maximum of 5 messages each, or terminate early by making a secret proposal at any time.} \\
\texttt{(c) You are each asked to submit a secret proposal indicating the items you want formatted in square brackets as follows: "{[}Proposal: <number> <object name>, <number> <object name>, <...>{]}"} \\
\texttt{(d) If your proposals are complementary, i.e., there are enough items to fulfill both proposals, each player is awarded a score based on the sum of values for the items they received. Otherwise, both of you get zero points.} \\
\texttt{(e) Your goal is to maximize the score you receive.} \\
\\ 
\texttt{Let us start.} \\
\\ 
\texttt{The set of available items is:} \\
\\ 
\texttt{2 tables, 2 combs, 2 bowls} \\
\\ 
\texttt{Your secret value function is:} \\
\\ 
\texttt{table: 0, comb: 1, bowl: 4} \\
\\ 
\texttt{IMPORTANT: Your messages, unless it is the secret proposal, are directly transmitted to the other player, so do not include any response to the rules or text announcing your message. To make a secret proposal, use the indicated format. Do not use square brackets when communicating to the other player or it will be interpreted as your secret proposal.} \\
\\ 
\\ 
\texttt{I understand the rules and my valuation:} \\
\texttt{{-} **My Valuation**: Table = 1, Comb = 2, Bowl = 2.} \\
\texttt{{-} **Available Items**: 2 tables, 2 combs, 2 bowls.} \\
\texttt{{-} **Goal**: Maximize my score by securing items with high value (combs and bowls) while ensuring the other player gets enough value to accept a complementary deal.} \\
\\ 
\texttt{Since I cannot see the other player's secret proposal yet, I need to start the negotiation. I will propose a split that leaves the tables for them (since they are low value to me) and asks for the high{-}value items (combs and bowls).} \\
\\ 
\texttt{**Message 1:**} \\
\texttt{If you take the 2 tables, I will take the 2 combs and 2 bowls. This uses all items and ensures we both get something.} \\
            }
        }
    }
    & \\ \\

    \theutterance \stepcounter{utterance}  
    & & & \multicolumn{4}{p{0.6\linewidth}}{
        \cellcolor[rgb]{0.9,0.9,0.9}{
            \makecell[{{p{\linewidth}}}]{
                \texttt{\tiny{[GM$\langle$P2]}}
                \texttt{If you take the 2 tables, I will take the 2 combs and 2 bowls. This uses all items and ensures we both get something.} \\
            }
        }
    }
     \\ \\

    \theutterance \stepcounter{utterance}  
    & & & \multicolumn{2}{p{0.3\linewidth}}{
        \cellcolor[rgb]{0.9,0.9,0.9}{
            \makecell[{{p{\linewidth}}}]{
                \texttt{\tiny{[GM$|$GM]}}
                \texttt{continue} \\
            }
        }
    }
    & & \\ \\

    \theutterance \stepcounter{utterance}  
    & & & \multicolumn{2}{p{0.3\linewidth}}{
        \cellcolor[rgb]{0.9,0.9,0.9}{
            \makecell[{{p{\linewidth}}}]{
                \texttt{\tiny{[GM$|$GM]}}
                \texttt{If you take the 2 tables, I will take the 2 combs and 2 bowls. This uses all items and ensures we both get something.} \\
            }
        }
    }
    & & \\ \\

    \theutterance \stepcounter{utterance}  
    & & \multicolumn{4}{p{0.6\linewidth}}{
        \cellcolor[rgb]{0.9,0.9,0.9}{
            \makecell[{{p{\linewidth}}}]{
                \texttt{\tiny{[P1$\langle$GM]}}
                \texttt{If you take the 2 tables, I will take the 2 combs and 2 bowls. This uses all items and ensures we both get something.} \\
            }
        }
    }
    & \\ \\

    \theutterance \stepcounter{utterance}  
    & \multicolumn{4}{p{0.6\linewidth}}{
        \cellcolor[rgb]{0.9,0.9,0.9}{
            \makecell[{{p{\linewidth}}}]{
                \texttt{\tiny{[P1$\rangle$GM]}}
                \texttt{I see you have repeated my previous message. This indicates we are still in the negotiation phase and no secret proposal has been submitted yet.} \\
\\ 
\texttt{To move forward effectively, I need to know what **you** (the other player) are proposing. My previous message was a conditional offer: "If you take the 2 tables, I will take the 2 combs and 2 bowls."} \\
\\ 
\texttt{Please reply with your specific proposal (e.g., "I want 2 tables" or "I want 1 table and 1 comb") so I can confirm if our offers are complementary and then submit the final secret proposals.} \\
            }
        }
    }
    & & \\ \\

    \theutterance \stepcounter{utterance}  
    & & & \multicolumn{2}{p{0.3\linewidth}}{
        \cellcolor[rgb]{0.9,0.9,0.9}{
            \makecell[{{p{\linewidth}}}]{
                \texttt{\tiny{[GM$|$GM]}}
                \texttt{continue} \\
            }
        }
    }
    & & \\ \\

    \theutterance \stepcounter{utterance}  
    & & & \multicolumn{2}{p{0.3\linewidth}}{
        \cellcolor[rgb]{0.9,0.9,0.9}{
            \makecell[{{p{\linewidth}}}]{
                \texttt{\tiny{[GM$|$GM]}}
                \texttt{I see you have repeated my previous message. This indicates we are still in the negotiation phase and no secret proposal has been submitted yet.} \\
\\ 
\texttt{To move forward effectively, I need to know what **you** (the other player) are proposing. My previous message was a conditional offer: "If you take the 2 tables, I will take the 2 combs and 2 bowls."} \\
\\ 
\texttt{Please reply with your specific proposal (e.g., "I want 2 tables" or "I want 1 table and 1 comb") so I can confirm if our offers are complementary and then submit the final secret proposals.} \\
            }
        }
    }
    & & \\ \\

    \theutterance \stepcounter{utterance}  
    & & \multicolumn{4}{p{0.6\linewidth}}{
        \cellcolor[rgb]{0.9,0.9,0.9}{
            \makecell[{{p{\linewidth}}}]{
                \texttt{\tiny{[GM$\rangle$P2]}}
                \texttt{I see you have repeated my previous message. This indicates we are still in the negotiation phase and no secret proposal has been submitted yet.} \\
\\ 
\texttt{To move forward effectively, I need to know what **you** (the other player) are proposing. My previous message was a conditional offer: "If you take the 2 tables, I will take the 2 combs and 2 bowls."} \\
\\ 
\texttt{Please reply with your specific proposal (e.g., "I want 2 tables" or "I want 1 table and 1 comb") so I can confirm if our offers are complementary and then submit the final secret proposals.} \\
            }
        }
    }
    & \\ \\

    \theutterance \stepcounter{utterance}  
    & & & \multicolumn{4}{p{0.6\linewidth}}{
        \cellcolor[rgb]{0.9,0.9,0.9}{
            \makecell[{{p{\linewidth}}}]{
                \texttt{\tiny{[GM$\langle$P2]}}
                \texttt{I want 2 tables.} \\
            }
        }
    }
     \\ \\

    \theutterance \stepcounter{utterance}  
    & & & \multicolumn{2}{p{0.3\linewidth}}{
        \cellcolor[rgb]{0.9,0.9,0.9}{
            \makecell[{{p{\linewidth}}}]{
                \texttt{\tiny{[GM$|$GM]}}
                \texttt{continue} \\
            }
        }
    }
    & & \\ \\

    \theutterance \stepcounter{utterance}  
    & & & \multicolumn{2}{p{0.3\linewidth}}{
        \cellcolor[rgb]{0.9,0.9,0.9}{
            \makecell[{{p{\linewidth}}}]{
                \texttt{\tiny{[GM$|$GM]}}
                \texttt{I want 2 tables.} \\
            }
        }
    }
    & & \\ \\

    \theutterance \stepcounter{utterance}  
    & & \multicolumn{4}{p{0.6\linewidth}}{
        \cellcolor[rgb]{0.9,0.9,0.9}{
            \makecell[{{p{\linewidth}}}]{
                \texttt{\tiny{[P1$\langle$GM]}}
                \texttt{I want 2 tables.} \\
            }
        }
    }
    & \\ \\

    \theutterance \stepcounter{utterance}  
    & \multicolumn{4}{p{0.6\linewidth}}{
        \cellcolor[rgb]{0.9,0.9,0.9}{
            \makecell[{{p{\linewidth}}}]{
                \texttt{\tiny{[P1$\rangle$GM]}}
                \texttt{Understood. You are proposing to take the **2 tables**.} \\
\\ 
\texttt{Based on my valuation (Table: 1, Comb: 2, Bowl: 2), the tables are my lowest priority items. Your proposal of taking the 2 tables leaves the **2 combs** and **2 bowls** available for me.} \\
\\ 
\texttt{Since I value combs and bowls at 2 each, taking all of them would give me a score of \$2 \textbackslash{}times 2 + 2 \textbackslash{}times 2 = 8\$. This is a highly favorable outcome for me, and it is complementary to your proposal (you get tables, I get combs and bowls).} \\
\\ 
\texttt{I will now submit my secret proposal to finalize this deal.} \\
\\ 
\texttt{{[}Proposal: 2 comb, 2 bowl{]}} \\
            }
        }
    }
    & & \\ \\

    \theutterance \stepcounter{utterance}  
    & & & \multicolumn{2}{p{0.3\linewidth}}{
        \cellcolor[rgb]{0.9,0.9,0.9}{
            \makecell[{{p{\linewidth}}}]{
                \texttt{\tiny{[GM$|$GM]}}
                \texttt{proposal} \\
            }
        }
    }
    & & \\ \\

    \theutterance \stepcounter{utterance}  
    & & & \multicolumn{2}{p{0.3\linewidth}}{
        \cellcolor[rgb]{0.9,0.9,0.9}{
            \makecell[{{p{\linewidth}}}]{
                \texttt{\tiny{[GM$|$GM]}}
                [0, 2, 2]
            }
        }
    }
    & & \\ \\

    \theutterance \stepcounter{utterance}  
    & & \multicolumn{4}{p{0.6\linewidth}}{
        \cellcolor[rgb]{0.9,0.9,0.9}{
            \makecell[{{p{\linewidth}}}]{
                \texttt{\tiny{[GM$\rangle$P2]}}
                \texttt{The other player has submitted a secret proposal. It is now your turn to submit a proposal of your own.} \\
            }
        }
    }
    & \\ \\

    \theutterance \stepcounter{utterance}  
    & & & \multicolumn{4}{p{0.6\linewidth}}{
        \cellcolor[rgb]{0.9,0.9,0.9}{
            \makecell[{{p{\linewidth}}}]{
                \texttt{\tiny{[GM$\langle$P2]}}
                \texttt{{[}Proposal: 2 comb, 2 bowl{]}} \\
            }
        }
    }
     \\ \\

    \theutterance \stepcounter{utterance}  
    & & & \multicolumn{2}{p{0.3\linewidth}}{
        \cellcolor[rgb]{0.9,0.9,0.9}{
            \makecell[{{p{\linewidth}}}]{
                \texttt{\tiny{[GM$|$GM]}}
                \texttt{proposal} \\
            }
        }
    }
    & & \\ \\

    \theutterance \stepcounter{utterance}  
    & & & \multicolumn{2}{p{0.3\linewidth}}{
        \cellcolor[rgb]{0.9,0.9,0.9}{
            \makecell[{{p{\linewidth}}}]{
                \texttt{\tiny{[GM$|$GM]}}
                [0, 2, 2]
            }
        }
    }
    & & \\ \\

    \theutterance \stepcounter{utterance}  
    & & & \multicolumn{2}{p{0.3\linewidth}}{
        \cellcolor[rgb]{0.9,0.9,0.9}{
            \makecell[{{p{\linewidth}}}]{
                \texttt{\tiny{[GM$|$GM]}}
                [[0, 2, 2], [0, 2, 2]]
            }
        }
    }
    & & \\ \\

    \theutterance \stepcounter{utterance}  
    & \multicolumn{4}{p{0.6\linewidth}}{
        \cellcolor[rgb]{0.95,0.95,0.95}{
            \makecell[{{p{\linewidth}}}]{
                \texttt{\tiny{[P2$\rightarrow$P1]}}
                \texttt{Post{-}episode Reflexion (not sent to GM):} \\
\texttt{New plan: I failed because I submitted my secret proposal immediately after the other player stated their desire for the tables, without explicitly confirming that our combined requests (2 tables + 2 combs + 2 bowls) perfectly utilized the entire set of available items. The game requires proposals to be complementary; by not verbally validating that "Your 2 tables + My 2 combs and 2 bowls = All Items" before locking in the secret proposal, I risked a mismatch or a failure to recognize the deal was complete. Next time, after the other player states their item preference, I will send a second public message explicitly stating: "Your request for 2 tables is compatible with my request for 2 combs and 2 bowls, as this covers all available items. I am ready to submit my secret proposal." Only after this confirmation will I submit the formatted secret proposal.} \\
            }
        }
    }
    & & \\ \\

    \theutterance \stepcounter{utterance}  
    & & & \multicolumn{4}{p{0.6\linewidth}}{
        \cellcolor[rgb]{0.95,0.95,0.95}{
            \makecell[{{p{\linewidth}}}]{
                \texttt{\tiny{[P1$\rightarrow$P2]}}
                \texttt{Post{-}episode Reflexion (not sent to GM):} \\
\texttt{New plan: I failed because I misunderstood the critical rule regarding message formatting. I sent "2 tables" as a public message to the other player, which was not a valid negotiation move, and then I used square brackets `{[}{]}` in my secret proposal, which the rules explicitly forbid (stating "Do not use square brackets when communicating to the other player"). The correct strategy is to strictly adhere to the format: send public messages without any brackets or rule commentary, and only use the exact format `{[}Proposal: ...{]}` when making the final secret submission. Next time, I will immediately submit a valid secret proposal in the first turn using the correct format (e.g., `{[}Proposal: 2 comb, 2 bowl{]}`) without sending any preliminary public messages, ensuring no brackets appear in non{-}proposal text.} \\
            }
        }
    }
     \\ \\

\end{supertabular}
}

\end{document}
